\documentclass[11pt,a4paper]{article}
\usepackage[margin=1in]{geometry}
\usepackage{adjustbox}
\usepackage{placeins}
\title{MOTIF: Preliminary Link-Prediction and Invariance Results}
\author{}
\date{}
\begin{document}
\maketitle

\section{IMDb}
\subsection{Movie--Year: Zero-Shot}
The released checkpoint is held fixed across graph variants.
% Requires adjustbox.
\begin{table}[htbp]
\centering
\small
\renewcommand{\arraystretch}{1.15}
\setlength{\tabcolsep}{5pt}
\begin{adjustbox}{max width=\linewidth}
\begin{tabular}{|l|l|l|l|l|l|l|}
\hline
\textbf{Dataset} & \textbf{Method} & \textbf{Variant} & \textbf{MRR} & \textbf{Hits@1} & \textbf{Hits@3} & \textbf{Hits@10} \\
\hline
IMDb & MOTIF & v1 & $0.1802 \pm 0.0082$ & $0.0624 \pm 0.0085$ & $0.1379 \pm 0.0120$ & $0.4580 \pm 0.0175$ \\
\hline
IMDb & MOTIF & v2 & $0.1901 \pm 0.0087$ & $0.0755 \pm 0.0093$ & $0.1487 \pm 0.0123$ & $0.4652 \pm 0.0175$ \\
\hline
IMDb & MOTIF & v3 & $0.1553 \pm 0.0068$ & $0.0396 \pm 0.0068$ & $0.1055 \pm 0.0107$ & $0.4436 \pm 0.0171$ \\
\hline
IMDb & MOTIF & v4 & $0.1984 \pm 0.0090$ & $0.0791 \pm 0.0094$ & $0.1667 \pm 0.0134$ & $0.4748 \pm 0.0174$ \\
\hline
\end{tabular}
\end{adjustbox}
\caption{Preliminary zero-shot MOTIF results on IMDb Movie--Year; 834 test queries. Tail-only link prediction over the original 20 candidates; score ties use average rank. Values are estimates $\pm$ bootstrap standard error (1,000 query resamples).}
\label{tab:motif-imdb-movie-year-link_prediction}
\end{table}

\FloatBarrier
% Requires adjustbox.
\begin{table}[htbp]
\centering
\small
\renewcommand{\arraystretch}{1.15}
\setlength{\tabcolsep}{5pt}
\begin{adjustbox}{max width=\linewidth}
\begin{tabular}{|l|l|l|l|l|}
\hline
\textbf{Method} & \textbf{Comparison} & \textbf{Agreement@1} & \textbf{Kendall-$\tau_b$} & \textbf{$\Delta$MRR} \\
\hline
MOTIF & v1 vs v2 & $0.2410 \pm 0.0146$ & $0.6166 \pm 0.0085$ & $0.0100 \pm 0.0100$ \\
\hline
MOTIF & v1 vs v3 & $0.3465 \pm 0.0160$ & $0.6748 \pm 0.0075$ & $-0.0249 \pm 0.0085$ \\
\hline
MOTIF & v1 vs v4 & $0.7806 \pm 0.0147$ & $0.9101 \pm 0.0038$ & $0.0182 \pm 0.0040$ \\
\hline
MOTIF & v2 vs v3 & $0.6379 \pm 0.0174$ & $0.9114 \pm 0.0043$ & $-0.0349 \pm 0.0057$ \\
\hline
MOTIF & v2 vs v4 & $0.2950 \pm 0.0153$ & $0.6672 \pm 0.0082$ & $0.0082 \pm 0.0101$ \\
\hline
MOTIF & v3 vs v4 & $0.2518 \pm 0.0151$ & $0.6433 \pm 0.0083$ & $0.0431 \pm 0.0095$ \\
\hline
\end{tabular}
\end{adjustbox}
\caption{Preliminary zero-shot MOTIF results on IMDb Movie--Year; 834 test queries. The checkpoint is fixed across variants. Both rankings use the same 20 candidates per query. Agreement@1 compares predictions; Kendall tau-b compares rankings. Delta MRR is the second variant minus the first. Values are estimates $\pm$ bootstrap standard error (1,000 query resamples).}
\label{tab:motif-imdb-movie-year-invariance}
\end{table}

\FloatBarrier

\clearpage
\subsection{Movie--Year: Fine-Tuning}
Each variant starts independently from the released checkpoint and uses 1,000 optimizer steps. Validation MRR selects each checkpoint. Training seed: 1024.

Hardware: NVIDIA H100 80GB HBM3 MIG 4g.40gb.

Recorded training time / peak allocated GPU memory: v1: 54.18 s / 0.770 GiB; v2: 52.02 s / 0.770 GiB; v3: 51.65 s / 0.770 GiB; v4: 51.72 s / 0.770 GiB.

Batch target-edge and inverse removals: 7,994 per training run (7.994 per update on average).

This is a single-seed preliminary experiment. Error bars are query-bootstrap standard errors. The guide's repeated-training noise floor and multi-seed mean/standard deviation are not yet available.


The prediction-performance table has 16 rows, one for each combination of
Fine-tuned on and Test graph. Matching labels give matched-variant results
(C-a); the other 12 rows give transfer results (C-b).
The invariance table labels two comparisons: C-a compares separately trained
models on their own variants; C-b compares test variants with one checkpoint
fixed, including the matched-variant reference. Each invariance row compares
two prediction outputs: six C-a comparisons and 24 C-b comparisons in total.
For example, checkpoint v1 with test graphs v1 vs v2 compares that same
v1-trained checkpoint's predictions on the two graphs.
% Requires adjustbox.
\begin{table}[htbp]
\centering
\small
\renewcommand{\arraystretch}{1.15}
\setlength{\tabcolsep}{5pt}
\begin{adjustbox}{max width=\linewidth}
\begin{tabular}{|l|l|l|l|l|l|l|}
\hline
\textbf{Regime} & \textbf{Fine-tuned on} & \textbf{Test graph} & \textbf{MRR} & \textbf{Hits@1} & \textbf{Hits@3} & \textbf{Hits@10} \\
\hline
C-a & v1 & v1 & $0.2028 \pm 0.0090$ & $0.0839 \pm 0.0094$ & $0.1787 \pm 0.0131$ & $0.4772 \pm 0.0173$ \\
\hline
C-b & v1 & v2 & $0.1875 \pm 0.0087$ & $0.0683 \pm 0.0092$ & $0.1583 \pm 0.0127$ & $0.4580 \pm 0.0175$ \\
\hline
C-b & v1 & v3 & $0.1724 \pm 0.0075$ & $0.0528 \pm 0.0076$ & $0.1403 \pm 0.0114$ & $0.4460 \pm 0.0176$ \\
\hline
C-b & v1 & v4 & $0.2225 \pm 0.0095$ & $0.0971 \pm 0.0100$ & $0.2062 \pm 0.0138$ & $0.4976 \pm 0.0177$ \\
\hline
C-b & v2 & v1 & $0.2183 \pm 0.0092$ & $0.0983 \pm 0.0097$ & $0.1954 \pm 0.0134$ & $0.4988 \pm 0.0172$ \\
\hline
C-a & v2 & v2 & $0.1823 \pm 0.0085$ & $0.0659 \pm 0.0087$ & $0.1427 \pm 0.0123$ & $0.4592 \pm 0.0179$ \\
\hline
C-b & v2 & v3 & $0.1551 \pm 0.0070$ & $0.0420 \pm 0.0070$ & $0.1091 \pm 0.0108$ & $0.4365 \pm 0.0174$ \\
\hline
C-b & v2 & v4 & $0.2311 \pm 0.0096$ & $0.1067 \pm 0.0103$ & $0.2146 \pm 0.0142$ & $0.5180 \pm 0.0174$ \\
\hline
C-b & v3 & v1 & $0.2092 \pm 0.0094$ & $0.0851 \pm 0.0097$ & $0.1930 \pm 0.0139$ & $0.5084 \pm 0.0173$ \\
\hline
C-b & v3 & v2 & $0.1913 \pm 0.0089$ & $0.0767 \pm 0.0094$ & $0.1511 \pm 0.0124$ & $0.4616 \pm 0.0174$ \\
\hline
C-a & v3 & v3 & $0.2110 \pm 0.0093$ & $0.0923 \pm 0.0100$ & $0.1703 \pm 0.0129$ & $0.5156 \pm 0.0176$ \\
\hline
C-b & v3 & v4 & $0.2006 \pm 0.0095$ & $0.0839 \pm 0.0102$ & $0.1655 \pm 0.0129$ & $0.4688 \pm 0.0172$ \\
\hline
C-b & v4 & v1 & $0.2140 \pm 0.0093$ & $0.0887 \pm 0.0098$ & $0.2026 \pm 0.0134$ & $0.4928 \pm 0.0173$ \\
\hline
C-b & v4 & v2 & $0.1922 \pm 0.0088$ & $0.0707 \pm 0.0090$ & $0.1691 \pm 0.0133$ & $0.4796 \pm 0.0173$ \\
\hline
C-b & v4 & v3 & $0.1722 \pm 0.0077$ & $0.0564 \pm 0.0079$ & $0.1331 \pm 0.0119$ & $0.4400 \pm 0.0170$ \\
\hline
C-a & v4 & v4 & $0.2253 \pm 0.0094$ & $0.1007 \pm 0.0100$ & $0.2098 \pm 0.0138$ & $0.5096 \pm 0.0175$ \\
\hline
\end{tabular}
\end{adjustbox}
\caption{Prediction performance and cross-representation transfer. Preliminary fine-tuned MOTIF on IMDb Movie--Year; seed 1024, 834 test queries. Each variant uses 1,000 optimizer steps; validation MRR selects its checkpoint. Each query ranks the correct year against the same 19 fixed negative years. Each row evaluates one trained checkpoint on one test graph. Matching train/test variants give matched-variant evaluation (C-a); other cells give cross-representation transfer (C-b). Score ties use average rank. Values are estimates $\pm$ query-bootstrap standard error (1,000 resamples), not variation across training seeds.}
\label{tab:motif-imdb-movie-year-matrix-transfer-performance}
\end{table}

\FloatBarrier
% Requires adjustbox.
\begin{table}[htbp]
\centering
\small
\renewcommand{\arraystretch}{1.15}
\setlength{\tabcolsep}{5pt}
\begin{adjustbox}{max width=\linewidth}
\begin{tabular}{|l|l|l|l|l|l|}
\hline
\textbf{Regime} & \textbf{Checkpoint trained on} & \textbf{Test graphs compared} & \textbf{Agreement@1} & \textbf{Kendall-$\tau_b$} & \textbf{$\Delta$MRR} \\
\hline
C-a & v1 and v2 (separate) & v1 vs v2 & $0.4652 \pm 0.0177$ & $0.3386 \pm 0.0067$ & $-0.0205 \pm 0.0091$ \\
\hline
C-a & v1 and v3 (separate) & v1 vs v3 & $0.0432 \pm 0.0070$ & $0.2528 \pm 0.0097$ & $0.0082 \pm 0.0130$ \\
\hline
C-a & v1 and v4 (separate) & v1 vs v4 & $0.4281 \pm 0.0181$ & $0.1890 \pm 0.0111$ & $0.0224 \pm 0.0099$ \\
\hline
C-a & v2 and v3 (separate) & v2 vs v3 & $0.1487 \pm 0.0124$ & $0.3181 \pm 0.0084$ & $0.0287 \pm 0.0107$ \\
\hline
C-a & v2 and v4 (separate) & v2 vs v4 & $0.3597 \pm 0.0174$ & $0.2976 \pm 0.0075$ & $0.0430 \pm 0.0099$ \\
\hline
C-a & v3 and v4 (separate) & v3 vs v4 & $0.0312 \pm 0.0062$ & $0.0277 \pm 0.0073$ & $0.0143 \pm 0.0139$ \\
\hline
C-b & v1 & v1 vs v2 & $0.1739 \pm 0.0133$ & $0.3876 \pm 0.0073$ & $-0.0153 \pm 0.0119$ \\
\hline
C-b & v1 & v1 vs v3 & $0.3177 \pm 0.0159$ & $0.4727 \pm 0.0067$ & $-0.0304 \pm 0.0104$ \\
\hline
C-b & v1 & v1 vs v4 & $0.6499 \pm 0.0162$ & $0.7616 \pm 0.0052$ & $0.0197 \pm 0.0060$ \\
\hline
C-b & v1 & v2 vs v3 & $0.3082 \pm 0.0156$ & $0.5590 \pm 0.0058$ & $-0.0151 \pm 0.0100$ \\
\hline
C-b & v1 & v2 vs v4 & $0.2818 \pm 0.0153$ & $0.4146 \pm 0.0076$ & $0.0350 \pm 0.0119$ \\
\hline
C-b & v1 & v3 vs v4 & $0.2866 \pm 0.0150$ & $0.3731 \pm 0.0074$ & $0.0501 \pm 0.0110$ \\
\hline
C-b & v2 & v1 vs v2 & $0.4341 \pm 0.0174$ & $0.5165 \pm 0.0064$ & $-0.0360 \pm 0.0088$ \\
\hline
C-b & v2 & v1 vs v3 & $0.3106 \pm 0.0164$ & $0.4115 \pm 0.0088$ & $-0.0632 \pm 0.0100$ \\
\hline
C-b & v2 & v1 vs v4 & $0.7806 \pm 0.0146$ & $0.8280 \pm 0.0048$ & $0.0128 \pm 0.0038$ \\
\hline
C-b & v2 & v2 vs v3 & $0.2818 \pm 0.0159$ & $0.6368 \pm 0.0061$ & $-0.0272 \pm 0.0080$ \\
\hline
C-b & v2 & v2 vs v4 & $0.3969 \pm 0.0177$ & $0.4918 \pm 0.0071$ & $0.0488 \pm 0.0095$ \\
\hline
C-b & v2 & v3 vs v4 & $0.1595 \pm 0.0127$ & $0.3032 \pm 0.0094$ & $0.0760 \pm 0.0110$ \\
\hline
C-b & v3 & v1 vs v2 & $0.1295 \pm 0.0115$ & $0.2466 \pm 0.0100$ & $-0.0179 \pm 0.0117$ \\
\hline
C-b & v3 & v1 vs v3 & $0.1199 \pm 0.0116$ & $0.3918 \pm 0.0102$ & $0.0018 \pm 0.0113$ \\
\hline
C-b & v3 & v1 vs v4 & $0.5815 \pm 0.0170$ & $0.8247 \pm 0.0050$ & $-0.0086 \pm 0.0058$ \\
\hline
C-b & v3 & v2 vs v3 & $0.4736 \pm 0.0174$ & $0.6977 \pm 0.0073$ & $0.0197 \pm 0.0072$ \\
\hline
C-b & v3 & v2 vs v4 & $0.1595 \pm 0.0120$ & $0.3052 \pm 0.0098$ & $0.0093 \pm 0.0115$ \\
\hline
C-b & v3 & v3 vs v4 & $0.0887 \pm 0.0097$ & $0.3003 \pm 0.0099$ & $-0.0104 \pm 0.0113$ \\
\hline
C-b & v4 & v1 vs v2 & $0.2602 \pm 0.0155$ & $0.1068 \pm 0.0070$ & $-0.0217 \pm 0.0105$ \\
\hline
C-b & v4 & v1 vs v3 & $0.4125 \pm 0.0169$ & $0.2969 \pm 0.0073$ & $-0.0418 \pm 0.0095$ \\
\hline
C-b & v4 & v1 vs v4 & $0.7206 \pm 0.0153$ & $0.7498 \pm 0.0054$ & $0.0113 \pm 0.0048$ \\
\hline
C-b & v4 & v2 vs v3 & $0.4700 \pm 0.0176$ & $0.5366 \pm 0.0066$ & $-0.0200 \pm 0.0067$ \\
\hline
C-b & v4 & v2 vs v4 & $0.2554 \pm 0.0149$ & $0.1105 \pm 0.0076$ & $0.0330 \pm 0.0113$ \\
\hline
C-b & v4 & v3 vs v4 & $0.2878 \pm 0.0156$ & $0.3154 \pm 0.0077$ & $0.0531 \pm 0.0101$ \\
\hline
\end{tabular}
\end{adjustbox}
\caption{Preliminary fine-tuned MOTIF on IMDb Movie--Year; seed 1024, 834 test queries. Each variant uses 1,000 optimizer steps; validation MRR selects its checkpoint. Each query ranks the correct year against the same 19 fixed negative years. C-a compares separately trained models, each on its own variant; weights and graphs both change. C-b groups hold the selected training checkpoint fixed, including its matched-variant reference. Agreement@1 compares predictions; Kendall tau-b compares within-query rankings. Delta MRR is the second test variant minus the first. Values are estimates $\pm$ query-bootstrap standard error (1,000 resamples), not variation across training seeds.}
\label{tab:motif-imdb-movie-year-matrix-invariance}
\end{table}

\FloatBarrier

\clearpage
\section{DBLP}
\subsection{Paper--Conference: Zero-Shot}
The released checkpoint is held fixed across graph variants.
% Requires adjustbox.
\begin{table}[htbp]
\centering
\small
\renewcommand{\arraystretch}{1.15}
\setlength{\tabcolsep}{5pt}
\begin{adjustbox}{max width=\linewidth}
\begin{tabular}{|l|l|l|l|l|l|l|}
\hline
\textbf{Dataset} & \textbf{Method} & \textbf{Variant} & \textbf{MRR} & \textbf{Hits@1} & \textbf{Hits@3} & \textbf{Hits@10} \\
\hline
DBLP & MOTIF & v1 & $0.5810 \pm 0.0067$ & $0.3886 \pm 0.0093$ & $0.7243 \pm 0.0082$ & $0.9370 \pm 0.0050$ \\
\hline
DBLP & MOTIF & v2 & $0.5696 \pm 0.0069$ & $0.3800 \pm 0.0094$ & $0.7052 \pm 0.0084$ & $0.9280 \pm 0.0052$ \\
\hline
DBLP & MOTIF & v3 & $0.5727 \pm 0.0068$ & $0.3830 \pm 0.0092$ & $0.7089 \pm 0.0083$ & $0.9351 \pm 0.0050$ \\
\hline
\end{tabular}
\end{adjustbox}
\caption{Preliminary zero-shot MOTIF results on DBLP Paper--Conference; 2,666 test queries. Tail-only link prediction over the original 20 candidates; score ties use average rank. Values are estimates $\pm$ bootstrap standard error (1,000 query resamples).}
\label{tab:motif-dblp-paper-conference-link_prediction}
\end{table}

\FloatBarrier
% Requires adjustbox.
\begin{table}[htbp]
\centering
\small
\renewcommand{\arraystretch}{1.15}
\setlength{\tabcolsep}{5pt}
\begin{adjustbox}{max width=\linewidth}
\begin{tabular}{|l|l|l|l|l|l|}
\hline
\textbf{Dataset} & \textbf{Method} & \textbf{Variant} & \textbf{Precision (macro)} & \textbf{Recall (macro)} & \textbf{Macro-F1} \\
\hline
DBLP & MOTIF & v1 & $0.3203 \pm 0.0140$ & $0.3077 \pm 0.0090$ & $0.2955 \pm 0.0094$ \\
\hline
DBLP & MOTIF & v2 & $0.3339 \pm 0.0161$ & $0.2907 \pm 0.0084$ & $0.2831 \pm 0.0096$ \\
\hline
DBLP & MOTIF & v3 & $0.3231 \pm 0.0152$ & $0.2938 \pm 0.0084$ & $0.2847 \pm 0.0093$ \\
\hline
\end{tabular}
\end{adjustbox}
\caption{Preliminary zero-shot MOTIF results on DBLP Paper--Conference; 2,666 test queries. Top-1 conference prediction; precision, recall and F1 are macro-averaged over all 20 conferences. Accuracy and Micro-F1 equal Hits@1 in the ranking table and are omitted. Top-score ties use the smallest global conference ID. Values are estimates $\pm$ bootstrap standard error (1,000 query resamples).}
\label{tab:motif-dblp-paper-conference-conference_prediction}
\end{table}

\FloatBarrier
% Requires booktabs and adjustbox.
\begin{table}[htbp]
\centering
\small
\begin{adjustbox}{max width=\linewidth}
\begin{tabular}{lrr}
\toprule
Comparison & Agreement@1 & Kendall $\tau_b$ \\
\midrule
v1 vs v2 & $0.8953 \pm 0.0057$ & $0.9506 \pm 0.0005$ \\
v1 vs v3 & $0.9070 \pm 0.0058$ & $0.9570 \pm 0.0005$ \\
v2 vs v3 & $0.9156 \pm 0.0055$ & $0.9591 \pm 0.0004$ \\
\bottomrule
\end{tabular}
\end{adjustbox}
\caption{MOTIF / DBLP / Paper-Conference (20-candidate representation sensitivity). Preliminary frozen zero-shot evaluation of physical variants on the test split; 2666 queries. Both columns use the same original 20 candidates per query. Agreement@1 compares predicted entity IDs; Kendall tau-b compares candidate rankings within each query and is then averaged over queries. Values are estimates plus/minus bootstrap standard error from 1000 query resamples.}
\label{tab:motif-dblp-paper-conference-invariance}
\end{table}

\FloatBarrier

\end{document}
